\subsection*{Observed Values Are Not in This Folder, and Do Not Need to Be}

The evaluation joins forecasts to observed fatalities on \texttt{month\_id} and
\texttt{priogrid\_id}. Those observations are not included here, and Sonja's ensemble-construction
tool does not need them to be: it fetches actuals itself from a viewser queryset named in its own
configuration, and also accepts an actuals dataframe passed in directly. They can equally be pulled
from the dataFactory.

Three things about that tool are worth flagging to whoever runs it, taken from reading its current
source rather than from memory.

\begin{itemize}
    \item \textbf{The file naming already matches.} It expects\\
          \path{predictions_calibration_{timestamp}_{ts}.parquet}\\
          with \texttt{ts} running \texttt{00}--\texttt{12}, which is exactly what is here.
    \item \textbf{Its shipped default configuration is country-month, not PRIO-GRID.} The default
          keys are \texttt{(country\_id, month\_id)} and the default actuals queryset is a
          country-level one. The evaluation code itself is indifferent to the level --- it merges
          on whatever keys the configuration names --- so pointing it at these models is a
          configuration change, to \texttt{(priogrid\_id, month\_id)} and a PRIO-GRID actuals
          queryset, not a code change.
    \item \textbf{It resolves prediction files inside the repository layout}, at
          \path{views-models/models/<model>/data/generated/}, rather than from an
          arbitrary folder. The files here are laid out one folder per model and will need placing
          accordingly, or the paths adjusting.
\end{itemize}

\subsection*{If a Research Assistant Runs It Instead}

The instruction is: \textit{iterate over the thirteen files per model, join each to the observed
values on} \texttt{month\_id} \textit{and} \texttt{priogrid\_id}\textit{, and compute point metrics
per origin so that horizon effects stay separable.} Parquet reads in one call ---
\texttt{pandas.read\_parquet} in Python, \texttt{arrow::read\_parquet} in R.

Two things to tell them explicitly, because both will otherwise be guessed wrong:

\begin{enumerate}
    \item forecasts are counts, not logs --- no back-transformation;
    \item there is no distribution --- no distributional scores.
\end{enumerate}
